% =====================================================
% MÉTADONNÉES DU RAPPORT D'ANALYSE
% =====================================================

\renewcommand{\reporttitle}{Analyse des performances des campagnes marketing}
\renewcommand{\reportsubtitle}{De PostgreSQL à Python : exploration, analyse et visualisation de 200\,005 campagnes marketing}
\renewcommand{\reporttype}{RAPPORT D'ANALYSE DE DONNÉES}

\renewcommand{\reportauthor}{Manassé Lumu Bamanoa}
\renewcommand{\studentid}{}
\renewcommand{\program}{Data Analytics / Data Analysis}
\renewcommand{\academicyear}{2026}
\renewcommand{\internshipperiod}{Année 2021 (données) --- rédaction 2026}

\renewcommand{\university}{Data Analytics}
\renewcommand{\department}{Analyse de données}
\renewcommand{\academicsupervisor}{}
\renewcommand{\supervisorposition}{}

\renewcommand{\company}{PostgreSQL, SQL, Python, Pandas, Matplotlib, Streamlit}
\renewcommand{\companydepartment}{Pipeline analytique}
\renewcommand{\companytutor}{}
\renewcommand{\tutorposition}{}

\renewcommand{\defensedate}{2026}

\newcommand{\projectname}{Marketing Campaigns Analysis}
\newcommand{\maintech}{Python}
\newcommand{\backendtech}{PostgreSQL}
\newcommand{\databasetech}{PostgreSQL}

\hypersetup{
    pdftitle={Analyse des performances des campagnes marketing},
    pdfauthor={Manassé Lumu Bamanoa},
    pdfsubject={Data Analytics / Data Analysis},
    pdfkeywords={Data Analytics, PostgreSQL, SQL, Python, Pandas, Streamlit}
}
